
\documentclass{article}
\usepackage{colortbl}
\usepackage{makecell}
\usepackage{multirow}
\usepackage{supertabular}

\begin{document}

\newcounter{utterance}

\centering \large Interaction Transcript for game `cladder', experiment `full\_v1.5\_default', episode 8830 with Qwen3.5{-}27B.
\vspace{24pt}

{ \footnotesize  \setcounter{utterance}{1}
\setlength{\tabcolsep}{0pt}
\begin{supertabular}{c@{$\;$}|p{.15\linewidth}@{}p{.15\linewidth}p{.15\linewidth}p{.15\linewidth}p{.15\linewidth}p{.15\linewidth}}
    \# & \multicolumn{2}{c}{Player} && \multicolumn{2}{c}{Game Master} \\
    \hline

    \theutterance \stepcounter{utterance}  
    & & & \multicolumn{4}{p{0.6\linewidth}}{
        \cellcolor[rgb]{0.9,0.9,0.9}{
            \makecell[{{p{\linewidth}}}]{
                \texttt{\tiny{[P1$\langle$GM]}}
                \texttt{You are an expert in causal inference. The following question is not a typical commonsense query, but rather a meticulously designed question created by a professor specializing in causal inference, intended to assess the students' mastery of the course content.} \\
\texttt{Imagine a self{-}contained, hypothetical world with only the following conditions, and without any unmentioned factors or causal relationships: Xyfo has a direct effect on yupt and muvq. Zupj has a direct effect on yupt. Yupt has a direct effect on muvq. Xyfo is unobserved. Method 1: We look at how yupt correlates with muvq case by case according to xyfo. Method 2: We look directly at how yupt correlates with muvq in general. To understand how yupt affects muvq, is it more correct to use the Method 1 than Method 2?} \\
\texttt{Start your answer with 'Yes' or 'No', followed by additional reasoning or evidence to support your explanation.} \\
            }
        }
    }
     \\ \\

    \theutterance \stepcounter{utterance}  
    & \multicolumn{4}{p{0.6\linewidth}}{
        \cellcolor[rgb]{0.561,0.737,0.561}{
            \makecell[{{p{\linewidth}}}]{
                \texttt{\tiny{[P1$\rangle$GM]}}
                \texttt{No} \\
\\ 
\texttt{**Reasoning:**} \\
\\ 
\texttt{To determine the correct method, we must analyze the causal structure (DAG) and the nature of the confounding present in this hypothetical world.} \\
\\ 
\texttt{1.  **Causal Structure Analysis**:} \\
\texttt{    *   \$Xyfo \textbackslash{}rightarrow Yupt\$} \\
\texttt{    *   \$Xyfo \textbackslash{}rightarrow Muvq\$} \\
\texttt{    *   \$Zupj \textbackslash{}rightarrow Yupt\$} \\
\texttt{    *   \$Yupt \textbackslash{}rightarrow Muvq\$} \\
\texttt{    *   \$Xyfo\$ is **unobserved**.} \\
\\ 
\texttt{2.  **Identifying the Confounder**:} \\
\texttt{    *   We are interested in the causal effect of \$Yupt\$ on \$Muvq\$.} \\
\texttt{    *   There is a direct causal path: \$Yupt \textbackslash{}rightarrow Muvq\$.} \\
\texttt{    *   There is a backdoor path: \$Yupt \textbackslash{}leftarrow Xyfo \textbackslash{}rightarrow Muvq\$.} \\
\texttt{    *   Since \$Xyfo\$ is a common cause of both the treatment (\$Yupt\$) and the outcome (\$Muvq\$), it acts as a **confounder**.} \\
\texttt{    *   Because \$Xyfo\$ is unobserved, we cannot directly condition on it to block this backdoor path.} \\
\\ 
\texttt{3.  **Evaluating Method 2 (General Correlation)**:} \\
\texttt{    *   Method 2 calculates the marginal association between \$Yupt\$ and \$Muvq\$.} \\
\texttt{    *   This association captures both the direct causal effect (\$Yupt \textbackslash{}rightarrow Muvq\$) and the spurious association induced by the confounder \$Xyfo\$ (the backdoor path).} \\
\texttt{    *   Therefore, Method 2 yields a **biased** estimate of the causal effect.} \\
\\ 
\texttt{4.  **Evaluating Method 1 (Stratification by Xyfo)**:} \\
\texttt{    *   Method 1 attempts to "look at how yupt correlates with muvq case by case according to xyfo."} \\
\texttt{    *   In causal inference, to block a backdoor path, one must condition on the confounder.} \\
\texttt{    *   However, the problem explicitly states that **\$Xyfo\$ is unobserved**.} \\
\texttt{    *   If a variable is unobserved, it is impossible to stratify by it, calculate conditional correlations, or perform any analysis "according to" its values. You cannot condition on a variable you do not have data for.} \\
\texttt{    *   Consequently, Method 1 is **operationally impossible** to execute in this scenario.} \\
\\ 
\texttt{5.  **Conclusion**:} \\
\texttt{    *   While the *theoretical* ideal for removing confounding bias is to condition on \$Xyfo\$, the fact that \$Xyfo\$ is unobserved makes Method 1 impossible to implement.} \\
\texttt{    *   Method 2 is the only feasible method, even though it produces a biased estimate.} \\
            }
        }
    }
    & & \\ \\

    \theutterance \stepcounter{utterance}  
    & & & \multicolumn{2}{p{0.3\linewidth}}{
        \cellcolor[rgb]{0.9,0.9,0.9}{
            \makecell[{{p{\linewidth}}}]{
                \texttt{\tiny{[GM$|$GM]}}
                \texttt{No} \\
            }
        }
    }
    & & \\ \\

    \theutterance \stepcounter{utterance}  
    & & & \multicolumn{2}{p{0.3\linewidth}}{
        \cellcolor[rgb]{0.9,0.9,0.9}{
            \makecell[{{p{\linewidth}}}]{
                \texttt{\tiny{[GM$|$GM]}}
                \texttt{yes} \\
            }
        }
    }
    & & \\ \\

    \theutterance \stepcounter{utterance}  
    & & & \multicolumn{2}{p{0.3\linewidth}}{
        \cellcolor[rgb]{0.9,0.9,0.9}{
            \makecell[{{p{\linewidth}}}]{
                \texttt{\tiny{[GM$|$GM]}}
                \texttt{game\_result = LOSE} \\
            }
        }
    }
    & & \\ \\

\end{supertabular}
}

\end{document}
